\input BaTeX.tex

\PapierGroesse(400pt,300pt)

\begin{document}

\parskip=5pt
\baselineskip=13pt

Wie bleibt ein Kryptosystem sicher?\par
Entweder m"ussen System und Schl"ussel geheim bleiben oder es braucht mehr Schl"ussel, als ein Kryptoanalytiker effizient ausprobieren kann.

1883 formulierte Auguste Kerckhoff (1835-1903) ein Prinzip für die Sicherheit von Kryptosystemen, das heute noch gültig ist:
\par\bigskip

\hbox to\hsize{\hfill\vbox{\hsize=300pt\bf
Die Sicherheit des Kryptosystems muss auschliesslich auf der Geheimhaltung des Schlüssels beruhen, nicht aber auf der Geheimhaltung des Systems.}\hfill}

\end{document}